\documentclass[tikz,border=5pt]{standalone}
\usepackage{ctex}
\usepackage{amsmath,amssymb}
\usetikzlibrary{arrows.meta,positioning}
% Shared visual language for LFS architecture diagrams.
\RequirePackage{../../mymath}

\definecolor{LFSBlue}{HTML}{2563EB}
\definecolor{LFSBlueSoft}{HTML}{DBEAFE}
\definecolor{LFSGreen}{HTML}{15803D}
\definecolor{LFSGreenSoft}{HTML}{DCFCE7}
\definecolor{LFSOrange}{HTML}{C2410C}
\definecolor{LFSOrangeSoft}{HTML}{FFEDD5}
\definecolor{LFSRed}{HTML}{B91C1C}
\definecolor{LFSRedSoft}{HTML}{FEE2E2}
\definecolor{LFSPurple}{HTML}{7E22CE}
\definecolor{LFSPurpleSoft}{HTML}{F3E8FF}
\definecolor{LFSGray}{HTML}{475569}
\definecolor{LFSGraySoft}{HTML}{F1F5F9}

\tikzset{
  lfs picture/.style={
    font=\small,
    node distance=8mm and 12mm,
    >=Latex,
    every path/.style={line cap=round, line join=round}
  },
  block/.style={
    draw=LFSBlue,
    fill=LFSBlueSoft,
    rounded corners=2pt,
    line width=0.8pt,
    align=center,
    inner xsep=7pt,
    inner ysep=5pt,
    minimum height=10mm
  },
  compute/.style={
    block,
    draw=LFSPurple,
    fill=LFSPurpleSoft
  },
  storage/.style={
    block,
    draw=LFSGreen,
    fill=LFSGreenSoft
  },
  interface/.style={
    block,
    draw=LFSOrange,
    fill=LFSOrangeSoft
  },
  risk/.style={
    block,
    draw=LFSRed,
    fill=LFSRedSoft
  },
  neutral/.style={
    block,
    draw=LFSGray,
    fill=LFSGraySoft
  },
  decision/.style={
    diamond,
    draw=LFSOrange,
    fill=LFSOrangeSoft,
    line width=0.8pt,
    align=center,
    inner sep=2pt,
    aspect=2
  },
  group/.style={
    draw=LFSGray,
    rounded corners=3pt,
    dashed,
    line width=0.7pt,
    inner sep=6pt
  },
  arrow/.style={
    -{Latex[length=2.4mm,width=1.6mm]},
    draw=LFSGray,
    line width=0.9pt
  },
  data arrow/.style={
    arrow,
    draw=LFSBlue
  },
  control arrow/.style={
    arrow,
    draw=LFSOrange,
    dashed
  },
  residual/.style={
    arrow,
    draw=LFSGreen,
    rounded corners=4pt
  },
  label text/.style={
    font=\scriptsize,
    text=LFSGray,
    fill=white,
    inner sep=1.5pt
  },
  title/.style={
    font=\bfseries,
    text=LFSGray,
    align=center
  }
}


\begin{document}
\begin{tikzpicture}[my picture, >=Latex]

% Panel 1: Loss Spike Curve
\begin{scope}[xshift=0cm]
  \node[font=\small\bfseries, text=MyGray] at (2.6, 4.2) {损失尖峰 (Loss Spike) 演进轨迹};

  \draw[->, line width=0.6pt] (0, 0) -- (5.6, 0) node[right, font=\tiny] {训练步数 $t$};
  \draw[->, line width=0.6pt] (0, 0) -- (0, 3.8) node[above, font=\tiny] {训练损失 $\mathcal{L}$};

  % Normal descent before spike
  \draw[line width=1.1pt, draw=MyBlue]
    (0.2, 3.2) to[out=-30, in=170] (2.2, 1.4);

  % Spike
  \draw[line width=1.2pt, draw=MyRed]
    (2.2, 1.4) -- (2.5, 3.5) node[above, font=\tiny\bfseries, text=MyRed] {异常尖峰 (Spike)};

  % Divergence branch (dashed red)
  \draw[line width=0.9pt, dashed, draw=MyRed]
    (2.5, 3.5) -- (4.8, 3.6) node[right, font=\tiny, text=MyRed] {未抑制 $\to$ 训练发散崩溃};

  % Recovery branch (solid green)
  \draw[line width=1.1pt, draw=MyGreen]
    (2.5, 3.5) to[out=-80, in=175] (3.2, 1.5) to[out=-5, in=180] (5.2, 1.1)
    node[right, font=\tiny, text=MyGreen] {自愈/回退 $\to$ 恢复正常收敛};

  \node[font=\scriptsize, text=MyGray, align=center] at (2.6, -0.8) {
    伴随表象：梯度范数突增数十倍、注意力权重极端极化、\\
    动态损失缩放因子连续下溢减半。
  };
\end{scope}

% Separator
\draw[dashed, draw=MyGray!40, line width=0.5pt] (6.0, -1.2) -- (6.0, 4.4);

% Panel 2: Diagnostics & Recovery Flow
\begin{scope}[xshift=6.6cm]
  \node[font=\small\bfseries, text=MyGray] at (2.8, 4.2) {尖峰处置与自愈机制决策树};

  \tikzset{
    box/.style={draw=MyBlue, fill=MyBlueSoft, rounded corners=2pt, font=\scriptsize, align=center, inner sep=3pt, minimum width=48mm},
    alert/.style={draw=MyRed, fill=MyRedSoft, rounded corners=2pt, font=\scriptsize\bfseries, align=center, inner sep=3pt, minimum width=48mm},
    act/.style={draw=MyGreen, fill=MyGreenSoft, rounded corners=2pt, font=\scriptsize, align=center, inner sep=3pt, minimum width=48mm}
  }

  \node[box] (n1) at (2.8, 3.2) {监测到损失或梯度范数突增超过阈值 $c$};
  \node[alert, below=4mm of n1] (n2) {检查是否出现非有限值 (NaN / Inf)};
  \node[act, below=4mm of n2] (n3) {触发跳步：清空梯度，不更新参数和矩状态};
  \node[act, below=4mm of n3] (n4) {动态损失缩放调小：$s \leftarrow s / 2$};
  \node[box, below=4mm of n4] (n5) {若连续数步未恢复：回退前一健康检查点\\并跳过该数据批次 (Corrupted Batch)};

  \draw[->, line width=0.7pt, draw=MyGray] (n1) -- (n2);
  \draw[->, line width=0.7pt, draw=MyRed] (n2) -- node[right, font=\tiny, text=MyRed] {是} (n3);
  \draw[->, line width=0.7pt, draw=MyGray] (n3) -- (n4);
  \draw[->, line width=0.7pt, draw=MyGray] (n4) -- (n5);

\end{scope}

\end{tikzpicture}
\end{document}
